% === thesis-writer: Discussion section (skeleton with placeholders) ===
% NOTE: Cross-references assume the Results section lives in results_section_draft.tex
% with labels tab:main_results, tab:split_strategy, fig:rmse_over_horizon,
% fig:qualitative_examples (see that file). Update labels here if they differ.

\section{Discussion}\label{sec:Discussion}

% TODO: confirm whether this thesis has a separate Limitations section/chapter.
% If yes, move the closing paragraph of \ref{sec:disc-limitations-note} there
% and drop the heading below to \subsection{Limitations and Future Work} only
% if no separate section exists.

\subsection{Generalisation across subjects}\label{sec:disc-generalisation}

The gap between the LOSO results (Table~\ref{tab:main_results}) and the
within-subject split results (Table~\ref{tab:split_strategy}) speaks directly to
how much each model relies on subject-specific gait signatures versus
generalisable biomechanical structure. Under the within-subject split, the
training set contains trials from the same individuals as the test set, so the
model is effectively calibrated to each subject's own gait during training;
under the LOSO split, by contrast, it must predict these subject-specific
characteristics for an individual it has never observed.

For the Transformer encoder, mean RMSE was \SI{7.03}{\degree} under the
within-subject split versus \SI{14.31}{\degree} under LOSO, roughly a twofold
increase in error. This indicates that the encoder is unable to find a
generalisable pattern of biomechanical structure that transfers across
individuals. The TCN baseline exhibits a closely comparable pattern: mean RMSE
rises from \SI{12.30}{\degree} under the within-subject split to
\SI{24.50}{\degree} under LOSO, also roughly a twofold increase
(\num{1.99}$\times$, versus \num{2.03}$\times$ for the Transformer). Because
the \emph{relative} size of this generalisation gap is essentially identical
for both architectures, MAE pre-training does not appear to specifically
reduce the cost of subject-independent generalisation; rather, the fact that a
model with no pre-training at all (TCN) incurs an equally large relative
penalty suggests that this generalisation cost is largely intrinsic to the
task and to the size and inter-subject variability of the dataset, rather
than to the presence or absence of self-supervised pre-training.

This contrasts with the broader self-supervised learning literature, where
masked-reconstruction pre-training is typically motivated by improved
representation transfer under distribution shift \cite{PLACEHOLDER}. % TODO: cite MAE / BERT-style pretraining-for-transfer papers
[PLACEHOLDER: 1-2 sentences connecting the per-subject RMSE spread in
Figure~\ref{fig:per_subject_rmse} to specific subjects — e.g.\ does the model
fail consistently on the same subjects regardless of architecture, suggesting an
outlier subject (unusual gait, sensor placement, or anthropometry) rather than a
modelling deficiency?]

\subsection{Sources and distribution of prediction error}\label{sec:disc-error}

Figure~\ref{fig:rmse_over_horizon} and the qualitative traces in
Figure~\ref{fig:qualitative_examples} indicate where, rather than simply how
much, the models err. For the Transformer encoder, per-step RMSE does not grow
monotonically across the forecast horizon: it rises from \SI{8.65}{\degree} at
the first step to a peak of approximately \SI{16.95}{\degree} around
\SI{1.0}{\second}--\SI{1.1}{\second} in, then plateaus and declines slightly
toward the end of the window (Section~\ref{sec:results-horizon}). Error is also
unevenly distributed across activities (Section~\ref{sec:results-activity}):
stair ascent is the hardest activity (\SI{17.49}{\degree}), level walking and
stair descent are intermediate (\SI{15.04}{\degree} and \SI{15.18}{\degree}),
and ramp ascent/descent are easiest (\SI{11.97}{\degree} and
\SI{11.91}{\degree}). The residual analysis further shows a mean residual of
\SI{4.53}{\degree} (prediction minus ground truth), indicating that the
encoder systematically over-predicts knee angle rather than erring with
zero-mean noise. [PLACEHOLDER: state whether this systematic over-prediction
is itself concentrated in particular activities or forecast positions, e.g.\
by computing the mean (not just RMS) residual per activity and per step.]
These patterns are consistent with several structural limitations of the
input signal:

\begin{itemize}
    \item \textbf{Single-sensor input.} Only a thigh-mounted IMU is used
    (Section~\ref{sec:Materials}); no shank- or foot-mounted sensor and no EMG
    signal is available. Knee angle is a function of both thigh and shank
    kinematics, so the shank's contribution must be inferred indirectly — a task
    that is likely harder during ballistic swing-phase motion than during the
    more constrained stance phase.
    \item \textbf{Calibration is range-only, not kinematic.} The per-subject
    min-max normalisation (Eq.~\ref{eq:per_subject_minmax}) calibrates signal
    \emph{amplitude} but provides no subject-specific kinematic calibration
    (e.g.\ thigh length, joint-centre offset), which may limit how precisely
    absolute joint angle can be recovered for an unseen subject.
    \item \textbf{Activity-transition ambiguity.} Near transitions between
    activity modes, the recent context window spans two different gait
    patterns, which may be out-of-distribution relative to any single training
    activity segment.
    \item \textbf{Sensor noise and placement variability.} Soft-tissue
    artefact and inter-session strap-placement differences are not explicitly
    modelled and may contribute high-frequency error independent of the
    activity being performed.
\end{itemize}

Stair ascent and level walking — the two highest-error activities — both
involve more pronounced swing-phase flexion than ramp ascent/descent,
consistent with the single-sensor-input limitation above: the thigh IMU alone
may underdetermine shank kinematics most severely during the largest, most
ballistic flexion movements. Stair ascent additionally involves frequent
transitions into and out of discrete steps, consistent with the
activity-transition-ambiguity limitation.

[PLACEHOLDER: 1-2 sentences stating whether the achieved RMSE/MAE
(Table~\ref{tab:main_results}) is within a range considered clinically or
functionally acceptable for prosthetic/exoskeleton control. Needs a citation for
the relevant acceptability threshold — e.g.\ \cite{PLACEHOLDER}.]
% TODO: find literature threshold for "acceptable" continuous knee angle RMSE in
% control applications before finalising this claim.

\subsection{Computational considerations for deployment}\label{sec:disc-compute}

The deployment context motivating this work — wearable, embedded control of a
prosthesis or exoskeleton (Section~\ref{sec:Materials}) — imposes constraints
on model size and inference latency in addition to accuracy. As shown in
Table~\ref{tab:compute_comparison} (Section~\ref{sec:results-compute}), TCN is
both far smaller and roughly \num{17}$\times$ faster than the Transformer
encoder at batch size 1.

This latency difference has a direct bearing on real-time viability. At a
sampling rate of \SI{100}{\hertz}, the \num{137}-step forecast horizon spans
\SI{1.37}{\second} into the future. The Transformer encoder's
\SI{16.9}{\milli\second} single-sample inference latency exceeds the duration
of the first two forecast steps (\SI{20}{\milli\second}) combined, meaning
that by the time a prediction becomes available, its earliest one to two
predicted samples already correspond to elapsed rather than future time and
are effectively unusable for control; the remaining \num{135} steps, however,
still constitute a useable forecast window. Despite this latency cost, the
Transformer's substantial accuracy advantage over TCN (Table~\ref{tab:main_results})
means that the loss of the first few forecast samples does not, on its own,
rule out the encoder's viability for this application.

\subsection{Relation to prior work}\label{sec:disc-related-work}

The use of masked-reconstruction self-supervised pre-training for a
Transformer-based time series encoder follows the broader trend of adapting
masked autoencoding \cite{PLACEHOLDER} % TODO: he2022mae or similar MAE citation
from vision to time series and biomedical signal domains
\cite{PLACEHOLDER}. % TODO: time-series MAE / PatchTST-style pretraining citation
[PLACEHOLDER: state whether the observed effect of MAE pre-training on LOSO
performance (Section~\ref{sec:disc-generalisation}) is consistent with results
reported in that literature for small-data, high-inter-subject-variance
biomedical settings, or whether it diverges, and offer a brief hypothesis for
any divergence — e.g.\ dataset size, masking ratio, or pretraining/fine-tuning
epoch budget.]

Continuous joint angle estimation from wearable IMUs has been studied with a
range of architectures, including recurrent networks, TCNs, and more recently
Transformer variants \cite{PLACEHOLDER}. % TODO: cite prior IMU-to-joint-angle
% regression papers, ideally including any other ENABL3S-based studies
[PLACEHOLDER: 2-3 sentences comparing the achieved RMSE/MAE in spirit
(not exact numbers, since experimental conditions such as sensor count, window
length, and activity set differ across studies) to representative values
reported in that literature, and noting where this study's single-thigh-IMU,
multi-activity, LOSO setting is more or less constrained than prior work.]
TimesNet's frequency-domain period detection \cite{PLACEHOLDER} and TCN's
dilated-convolution receptive field \cite{PLACEHOLDER} represent two established
alternatives to attention-based sequence modelling for this task; [PLACEHOLDER:
1-2 sentences situating the Transformer encoder's performance relative to these
architectural families, beyond the direct numeric comparison already given in
Table~\ref{tab:main_results}].